\subsection{ORB-SLAM3 Integration}
\label{sec:orb_slam3_integration}

The VBEE library is designed to be easily integrated into existing KV-SLAM systems. This section covers the integration of libVBEE into ORB-SLAM3. As discussed in section \ref{sec:orb_slam3}, ORB-SLAM3 was selected due to its wide adoption, high performance, and inclusion of several SLAM extensions which can benefit from outdated map point culling.

The original intention was to perform VBEE operations on every image frame in real-time as it passed through the SLAM system. However, these operations are time-consuming and computationally expensive, which led to a reduction in frame throughput and localization accuracy. Instead, the decision was made to operate on keyframes rather than frames, and to move the VBEE operations to the ends of episodic runs, acting as a post-processing step rather than a real-time update.

\subsubsection{Map Point}

The Map Point class is modified to include a VBEE instance, allowing for observability based existence estimation on a per-map point basis. The Map Point class includes a merge method which is called when the SLAM system's optimization operations identify that two map points likely represent the same visual feature. This method is modified to call the VBEE merge method to combine the VBEE instances of the two map points. 

\subsubsection{VBEE Post-Processing Operations}

Run ends are determined by matching the current frame's timestamp to a timestamp in the splits file provided by the episodic datasets. Figure \ref{fig:damping_coefficient_selection} shows the pipeline of VBEE operations which take place at the ends of runs.

At the end of the run, several iterations of Global Bundle Adjustment are performed to minimize the global error of the map. This operation repositions keyframes and merges map points, which give the VBEE operations the greatest chance of success. Each keyframe from the most recent run is then processed in order of collection. For each keyframe, two groups of map points are generated: positive observations and negative observations. Positive observations are map points that the SLAM system observed within the keyframe, while negative observations are map points that could be projected onto the camera's sensor based on the keyframe's pose, but were not observed within the keyframe. For each positively observed map point, an update is submitted to its VBEE instance with a viewpoint of the keyframe's position, and a seen status of 1. Similarly, for each negatively observed map point, an update is submitted with a seen status of 0.

After all new keyframes have been processed, map point culling begins. For each map point in the map, if its probability of existence is found to be less than the $P(E)$ threshold, the map point is marked bad. Once marked bad, ORB-SLAM3 will perform the necessary operations to remove the map point from the map, and clean up any pointers to its instance.

\subsubsection{RANSAC Integration}

RANSAC is used for relocalization and loop closing in ORB-SLAM3. The base implementation selects map points with a uniform probability of selection. The integration of VBEE into the system allows the selection to instead be weighted by the map point's probability of being seen $P(S)$ from the given viewpoint. The conjecture is that by weighting for selection of map points with a high likelihood of being seen, the number of iterations required to fit a model to the data will be minimized.

For each map point, $P(S)$ is calculated as follows:
$$
P(S)=P(S|E)P(E)+P(S|\neg E)P(\neg E)
$$
where $P(S|E)$ is the observability model's estimate from the current position, and $P(S|\neg E)$ is the false positive rate parameter. A selection is made from the pool of map points, with each map point's probability of selection proportional to its probability of being seen. Instrumentation is added to the loop closing and relocalization methods to count the number of iterations the RANSAC process takes to find a solution. 